\section{Sparse two-meridian seed search with complete candidate pruning}
\label{sec:sparse-seeds}

Report 47 contains a second useful native optimization besides its orbit
certificates: a faster search for the two-meridian derivations of
Section~\ref{sec:two-meridian}. We integrate its reusable closure state,
structurally restricted candidate pairs and proper-closure pruning.
The representation arithmetic, source-bound certificate format and
search-independent verifier are unchanged. The delivered frozen baseline is
byte-identical to the previously maintained module, so the comparison does
not depend on reconstructing an approximate historical implementation.

\subsection{The implication system and its checked sparse regime}

A signed Wirtinger crossing $(o,u,v,\epsilon)$ supplies implications
$\{o,u\}\Rightarrow v$ and $\{o,v\}\Rightarrow u$. Coincident dependencies
count once. For seeds $S$, let $\operatorname{cl}(S)$ be the least set
containing $S$ and closed under these rules. A successful one- or two-seed
closure supplies an acyclic original-crossing derivation of all meridians.
The unchanged exact arithmetic then imposes every original relation and
returns an unknot or nontrivial-knot verdict. Exhaustive closure failure
still says only that this diagram lacks the required derivation; it remains
\texttt{INCONCLUSIVE} for recognition.

The sparse regime requires every singleton closure to be trivial. This is
equivalent to absence of a productive unary rule: a rule with one distinct
dependency and a different target. A productive unary rule immediately grows
its singleton. Conversely, without one, no first step can grow any singleton.
The implementation checks the original rule list before using sparse pairs.
It does not assume that every knot diagram satisfies this condition.

Suppose there are at least three arcs and a pair $S$ generates all arcs.
Some rule must derive the first arc outside $S$. Its prerequisites are in
$S$. They cannot consist of one arc in the checked regime, so they are
exactly the two elements of $S$, and its target is outside $S$.
Consequently every generating pair occurs among the distinct prerequisite
pairs of productive binary rules. With $r$ directed rules, at most $r$
pairs need examination. A knot diagram with $n$ crossings has $O(n)$ arcs
and at most $2n$ such rules, in place of the former quadratic pair set.
On two arcs, the sole pair is retained explicitly: it already contains the
whole universe and requires no first rule. Zero and one arc are handled
by the earlier empty/singleton logic.

The productive-unary guard is necessary. For example, crossings
$(0,0,1)$ and $(1,2,3)$ on four abstract arcs imply $0\Rightarrow1$ and
$\{1,2\}\Rightarrow3$. The pair $\{0,2\}$ generates all four arcs, but
it is not a productive binary prerequisite pair. The apparent candidates
$\{1,2\}$ and $\{1,3\}$ cannot derive arc zero. The maintained algorithm
therefore retains exhaustive pair enumeration whenever a productive unary
rule exists. The example concerns the implication logic, not a claim that
these two crossings are a valid four-arc knot diagram.

\subsection{Reuse closure arrays without changing derivation order}

The old engine allocates dependency counters for every attempted seed set,
paying $\Theta(v+r)$ initialization even if the trial reaches only one or
two of its $v$ arcs. The new engine allocates two arrays once. Advancing a
generation number makes older entries logically absent. A known-generation
stamp records enqueueing; a processed-generation stamp records removal from
the queue. The reached list is rebuilt only from actually visited vertices.

A rule fires when all its distinct dependencies have processed stamps for
the current generation and its target lacks a current known stamp. The
processed distinction is essential for preserving the exact previous trace.
In the old counter algorithm, the remaining counter is precisely the number
of dependencies not yet removed from the queue. Thus the two firing
conditions agree after every dequeue. Induction on queue events gives the
same enqueues, rule choices and derivation order, including coincident
dependencies. Testing only the known stamps would still compute closure,
but could change the timing of rule firings and the certificate trace.

One attempt pays for its dequeued vertices and their rule incidences, not
for untouched vertices. Both stamps are needed because a vertex may have
been enqueued but not yet processed when another dependency is examined.
Generation counters have $O(\log(n+2))$ bits over the at-most-quadratic number
of attempts; no array of all possible pairs is allocated.

\subsection{Prune pairs inside failed closed sets}

If a failed trial exhausts its queue in a proper closed set $C$, then every
seed pair $S'\subseteq C$ also fails:
\[
 \operatorname{cl}(S')\subseteq\operatorname{cl}(C)=C\ne V.
\]
The producer stores candidate pairs as a sparse undirected graph. After a
failed trial, it examines candidate edges incident to reached vertices and
marks later edges whose other endpoint is also reached. It skips those edges
on subsequent enumeration. Only proper completed closures justify pruning;
a resource-interrupted prefix proves no failed candidate.

Candidates remain in lexicographic order. Every omitted or skipped pair is
proved unsuccessful, so the first successful pair remains the same as in the
old exhaustive algorithm. Singleton attempts are also retained in their
original order. Combined with identical queue traces, this preserves the
entire successful certificate, including its exact arithmetic and source
binding. The number of charged attempts now counts actually examined seeds;
proved failures need not consume a full closure computation. Under a fixed
budget more of the class may be exhausted, but no monotonic improvement for
every tiny resource allowance is promised: the new graph has a preparation
cost of its own.

\subsection{Search bounds and their limits}

Let $v$ be the number of arcs, $r$ the number of rules and $\mathcal A$ the
seed sets actually examined. For a reached closure $C_S$, let $d_R(x)$ and
$d_P(x)$ be its vertices' rule-incidence and candidate-graph degrees.
A coarse output-sensitive bound is
\[
 O\!\left(v+r\log(r+1)+
 \sum_{S\in\mathcal A}\sum_{x\in C_S}
 (1+d_R(x)+d_P(x))\right)
\]
word operations, with $O(v+r)$ working words. The graph term is needed only
for failed candidate-pair closures. This formula uses ordinary amortized
set lookup for pair deduplication; sorting and adjacent deduplication offer
a deterministic comparison alternative. Even quadratic worst-case hash
preparation leaves the following quadratic Wirtinger bound intact.

There are $v$ singleton trials and at most $r$ candidate-pair trials, each
visiting at most $O(v+r)$ incidences. For a validated knot diagram,
$v,r=O(n)$, giving $O(n^2)$ search word operations or
$O(n^2\log(n+2))$ elementary bit operations for indices. This replaces the
old cubic closure-search bound on the checked regime. The unrestricted
productive-unary fallback retains $O(v^2(v+r))$ search work, although its
individual attempts still benefit from reused state. Memory remains linear
in the rule system rather than in the number of potential pairs.

This theorem concerns finding the derivation. Successful representation
arithmetic, gcd/parity checks and mandatory certificate replay remain as in
Section~\ref{sec:two-meridian}; their deliberately conservative
$O(n^4\log(n+2))$ complete-stage bound is not improved merely by changing
closure enumeration. The public stage remains opt-in in the ordinary
pipeline, after invariant filters, with the same time, work and attempt caps.
Earlier filter costs are outside the standalone bound. Exhaustive failure
proves no knot-type result, and no cheap conversion of every knot diagram
into the two-seed class has been established.

\subsection{Native validation}

All six delivered source, test, corpus, benchmark and documentation files
were verified against the archived report-47 manifest at intake. The maintained
research README now distinguishes fresh native evidence from the delivery's
historical measurements. The previous maintained module and the delivered frozen
baseline both have SHA256 beginning \texttt{5eb36230e9b9}. The optimized
module is integrated without changing the delivered search implementation.
Fourteen focused tests pass in 2.309 seconds. They include all 4,096
four-arc systems of distinct crossing triples and all 512 three-arc systems
allowing productive unary rules. For every one- and two-seed set, a separate
fixed-point sweep agrees with the stamped engine, whose trace also agrees
exactly with the old counter engine. The tests compare the first successful
pair before and after pruning. Four hundred additional random degenerate
systems cover coincident labels, empty rules and repeated generations.

The twenty frozen actual-diagram fixtures retain exactly the same successful
certificates. The suite exercises proper-closure pruning, the planar
productive-unary fallback, search-independent replay, exact shared-work
boundaries and global cancellation. These checks preserve the distinction
between exhausting a seed class and completing knot recognition.

The full maintained suite passes 915 tests with Regina in 209.658 seconds.
A further seeded audit compares 240 held-out actual planar braid closures
against the unchanged frozen module. Verdicts and successful certificates
agree exactly: 102 unknots, 86 nontrivial knots and 52 exhaustive two-seed
failures. There are 153 sparse-regime inputs and 87 productive-unary fallback
inputs. Both old and new verifiers accept each conclusive certificate.
This comparison checks the search integration; the pre-existing geometric
and quaternion tests supply the independent arithmetic controls.

The same audit constructs disconnected three-vertex implication blocks to
check scaling of unsuccessful closures. With 5,000 blocks there are 15,000
vertices and 10,000 directed rules. The candidate graph has 10,000 edges,
and closed-set pruning skips 5,000 of them. There are 20,000 actual trials,
30,000 vertex dequeues, 40,000 visited rule incidences and 20,000 pruning
incidences, versus 112,492,500 potential unordered pairs. The corresponding
counts scale exactly with the number of blocks in four audited sizes.
Preparation's charged sorting term is retained. These are synthetic
implication systems, not asserted knot diagrams. The 0.508-second audit,
all signed planar inputs, the unary guard counterexample and exact source
hashes are in \path{data/sparse-seeds-native-audit.json}.

\subsection{Standalone and whole-pipeline measurements}

The benchmark retains the delivery's fixed twenty-diagram corpus and seed
2610081249, but reruns it against the current maintained pipeline. Standalone
arms are the frozen baseline, an identical baseline control, stamped closure
without candidate restriction, candidates without closed-set pruning, and
the complete optimization. Each has one shuffled warmup and nine measured
rounds. Every timed call includes fresh PD validation, seed search, successful
representation arithmetic and replay. An exhaustive unsuccessful seed search
counts as a completed class query but stays inconclusive for recognition.

\begin{center}
\begin{tabular}{@{}lrrrrr@{}}
\toprule
Input & old ms & stamped ms & pairs ms & full ms & ratio \\
\midrule
survivor-00 & 1.303 & 1.007 & 0.935 & 0.923 & 1.410 \\
survivor-01 & 1.244 & 0.754 & 0.521 & 0.424 & 2.917 \\
survivor-04 & 0.734 & 0.678 & 0.718 & 0.703 & 1.056 \\
survivor-09 & 0.628 & 0.593 & 0.625 & 0.631 & 1.002 \\
gordian & 456.551 & 40.373 & 4.837 & 4.581 & 100.120 \\
trefoil & 0.237 & 0.243 & 0.251 & 0.249 & 0.951 \\
conway & 0.697 & 0.501 & 0.401 & 0.369 & 1.882 \\
kinoshita\_terasaka & 0.710 & 0.508 & 0.409 & 0.435 & 1.791 \\
\bottomrule
\end{tabular}
\end{center}


Whole-pipeline measurements separately use the ordinary configuration with
a three-second budget and the compressed-group configuration with an
18-second total budget and 15-second group allowance. Each compares the
old stage, its identical control, the new stage and disabling the stage,
with one shuffled warmup and five measured rounds. Local two-meridian caps
remain 50 ms, two million work units and 10,000 attempted seeds. A pipeline
\texttt{UNKNOWN} remains a censored sample; it is not assigned a completion
time or used in paired completed-time ratios. All raw samples, actual stage
invocations, proof hashes and loaded-dependency hashes are retained in
\path{../fast/results/two_meridian_search_native_20261008.json}.

\begin{center}
\begin{tabular}{@{}lrrrr@{}}
\toprule
Input & ordinary old/new ms & ratio & group old/new ms & ratio \\
\midrule
survivor-00 & 3.003/2.541 & 1.184 & 3.004/2.624 & 1.158 \\
survivor-01 & 14.347/13.218 & 1.085 & 16.406/15.616 & 1.063 \\
survivor-04 & 1.742/1.731 & 1.005 & 1.749/1.728 & 1.010 \\
survivor-09 & 1.911/2.106 & 0.974 & 1.607/1.594 & 1.005 \\
gordian & --/-- & -- & 1683.188/1692.476 & 1.037 \\
trefoil & 0.055/0.056 & 0.998 & 0.054/0.054 & 1.004 \\
conway & 1.023/1.024 & 1.001 & 1.019/1.020 & 1.001 \\
kinoshita\_terasaka & 1.052/1.047 & 1.005 & 1.048/1.058 & 0.999 \\
\bottomrule
\end{tabular}
\end{center}


All 900 measured standalone calls finish their class query. On Gordian,
the old median is 456.551 ms and the optimized median is 4.581 ms; the
median paired ratio is 100.120. Stamped state alone takes 40.373 ms and
restricting to prerequisite pairs takes 4.837 ms. The number of attempted
seed sets falls from 10,011 to 274, and charged work from 4,347,002 to
12,251. With the unchanged two-million-work cap, the old search stops
before exhaustion whereas the new search completely excludes its seed
class. Both answers remain inconclusive for knot recognition. This is a
search improvement on a hard negative instance for the derivation class,
not a new obstruction to unknottedness.

The smaller Conway and Kinoshita--Terasaka fixtures have paired standalone
ratios 1.882 and 1.791. Other fixtures gain little: survivor-09 has ratio
1.002, and the trefoil has ratio 0.951. Metadata construction can cost more
than it saves when a tiny successful derivation is found immediately.
The identical-baseline controls range from 0.975 to 1.038 across the
standalone corpus. These measurements support the large search improvement
and its mechanism; they do not support a universal constant-factor gain.

There are 800 measured whole-pipeline calls, of which 780 complete. All
20 censored calls are ordinary-configuration Gordian runs, across the four
arms. The group configuration completes Gordian through the existing
cyclic-group backend. Its old/new medians are 1,683.188/1,692.476 ms,
while its median of paired old/new ratios is 1.037. The latter statistic
need not equal the ratio of separate medians. With only five rounds and
identical-control ratios ranging from 0.909 to 1.035 across the group corpus,
this is weak evidence for a whole-pipeline gain. Disabling the seed stage
has Gordian median 1,653.628 ms. The optimized stage quickly finishes its
unsuccessful class query; most elapsed time belongs to other work.

Survivor-00 improves from 3.003 to 2.541 ms in the ordinary pipeline
(paired ratio 1.184), while several invariant-filtered fixtures never reach
the seed stage and consequently show no systematic change. All timings,
including slower arms and censored samples, are retained. The measured
loaded-source hashes are unchanged throughout the run. The public default
therefore remains disabled: the evidence justifies integrating the faster
search for callers that enable this class test, without yet establishing
an automatic dispatch policy that improves the broader input distribution.

\paragraph{Remaining theoretical gap.}
The new bound is quadratic for closure discovery in its checked regime,
with an exact fallback outside it. It neither bounds the number of
Reidemeister moves required to enter the two-seed class nor supplies a
quasi-polynomial complete recognizer for arbitrary planar diagrams.
Future adaptive scheduling can use reached-closure sizes and charged work
as progress signals, but any early departure must remain inconclusive and
leave a complete fallback available. Such scheduling also requires measured
stage costs rather than extrapolating the standalone 100-fold ratio to
whole-pipeline performance.
